\section{Related Work}\label{sec:related}

\subsection{Spatiotemporal traffic forecasting}
\prev{Graph-based traffic forecasting models differ in how they combine spatial aggregation and temporal prediction. DCRNN uses bidirectional diffusion within a recurrent encoder--decoder \cite{li2018dcrnn}. STGCN uses a fully convolutional design \cite{yu2018stgcn}, while T-GCN combines graph convolution with gated recurrent units \cite{zhao2020tgcn}. Graph WaveNet combines diffusion-based spatial aggregation with dilated temporal convolutions and adaptive adjacency \cite{wu2019graphwavenet}. ASTGCN weights spatial and temporal dependencies through attention \cite{guo2019astgcn}, and GMAN uses spatial and temporal attention blocks with an encoder--decoder design \cite{zheng2020gman}. These methods establish different forecasting backbones; our study evaluates graph augmentation while holding the backbone fixed within each comparison.}

\prev{The contribution of the architecture should also be distinguished from the representation of sensor and time identities. STID shows that identity embeddings with a multilayer perceptron can form a competitive forecasting baseline \cite{shao2022stid}. STAEformer combines spatiotemporal adaptive embeddings with temporal and spatial Transformer layers \cite{liu2023staeformer}. We therefore treat improved forecasting accuracy as evidence about the evaluated predictor and graph together. It does not, by itself, identify an information-propagation mechanism.}

\subsection{Learning predictive sensor dependencies}
\prev{Several forecasters learn relationships that are not specified by the road graph. Graph WaveNet learns an adaptive dependency matrix from sensor embeddings \cite{wu2019graphwavenet}. MTGNN learns directed dependencies jointly with spatial and temporal processing \cite{wu2020mtgnn}. GTS optimizes a distribution over discrete graphs \cite{shang2021gts}, while AGCRN combines adaptive graph generation with node-specific parameter learning \cite{bai2020agcrn}. STFGNN constructs temporal relationships from similarities between traffic series and combines them with spatial structure \cite{li2021stfgnn}. These approaches demonstrate that informative sensor relationships can extend beyond geographical adjacency.}

\prev{Our selectors also use traffic observations, but construct a fixed set of links before training the neural forecaster. The original road graph is retained, and the main sparse comparisons use the same link budget. Graph WaveNet's adaptive adjacency is disabled in these comparisons so that the supplied graph determines its spatial supports. This design evaluates explicit link-selection rules within a common forecasting model; it does not compare against the full range of learned-graph architectures.}

\subsection{Over-squashing and structural rewiring}
\prev{Alon and Yahav identify the compression of growing graph neighbourhoods into fixed-size representations as a limitation of message passing \cite{alon2021bottleneck}. Topping et al. connect bottlenecks to graph curvature and propose curvature-based rewiring \cite{topping2022curvature}. Di Giovanni et al. analyse how width, depth and topology affect sensitivity to distant inputs \cite{digiovanni2023oversquashing}. Black et al. relate information propagation to effective resistance and propose Greedy Total Resistance (GTR) rewiring \cite{black2023resistance}. Effective resistance is also connected to random-walk commute times \cite{chandra1996resistance}. These results motivate resistance as a structural criterion, while leaving its forecasting value to be evaluated empirically.}

\prev{Other methods balance connectivity against different graph properties. FoSR adds edges using spectral expansion and distinguishes original from added connections through a relational architecture \cite{karhadkar2023fosr}. DiffWire uses differentiable layers motivated by commute time and spectral-gap optimization \cite{arnaiz2022diffwire}. Expander Graph Propagation uses a sparse expander as a communication structure \cite{deac2022egp}. LASER explicitly considers locality and sparsity alongside over-squashing \cite{barbero2024laser}. DRew changes multi-hop connectivity by layer and introduces distance-dependent delays \cite{gutteridge2023drew}. Graph Diffusion Convolution provides a broader graph-preprocessing perspective by sparsifying a generalized diffusion operator \cite{gasteiger2019diffusion}; it is relevant to changing communication neighbourhoods without being a traffic-specific over-squashing diagnostic.}

\prev{These methods do not use the same selection criteria or evaluation tasks. In particular, FoSR reports graph-classification experiments \cite{karhadkar2023fosr}, whereas our task is multi-horizon traffic-speed regression. The Long Range Graph Benchmark provides tasks for evaluating long-range dependencies \cite{dwivedi2022lrgb}. Such evaluations motivate testing distant information use, but their results cannot be assumed to transfer to low-speed traffic forecasting. In our experiments, GTR is a structural control; the other rewiring methods discussed here provide context and were not evaluated as forecasting baselines.}

\subsection{Spatiotemporal propagation, attention and virtual nodes}
\prev{Spatial graph structure is only one part of information propagation in a spatiotemporal model. Marisca et al. analyse sensitivity on joint spatial and temporal paths and compare temporal and spatial processing arrangements \cite{marisca2025stoversquashing}. This distinction matters because a useful source must be accessible at an informative time as well as at a useful location. Our lead--lag criterion addresses temporal ordering in the observed speeds, while our ridge screen evaluates one source summary for a 60-minute target. Neither criterion is a direct measurement of the neural model's spatiotemporal sensitivity.}

\prev{Global communication is also available through attention and hubs. Graphormer incorporates structural information into a Transformer \cite{ying2021graphormer}, and GraphGPS combines structural or positional encodings, local message passing and global attention \cite{rampasek2022gps}. For traffic, PDFormer models propagation delay and spatial relationships through attention \cite{jiang2023pdformer}. Cai et al. establish connections between message passing with virtual nodes and graph Transformers \cite{cai2023mpnnvn}. Cao et al. incorporate virtual nodes into STGCN using adaptive and semi-adaptive adjacency \cite{cao2025virtual}. Our virtual-node control follows the hub-connectivity idea, using one fixed global hub with neutral input. It is not a reproduction of Cao et al.'s learned semi-adaptive model, and its edge budget is larger than that of the sparse selectors.}

\subsection{Predictive screening and Granger-style construction}
\prev{Granger-style analysis asks whether the past of one series improves prediction of another after conditioning on the information already included \cite{granger1969}. STGC-GNNs uses spatiotemporal alignment and Granger testing to construct dependencies for traffic prediction \cite{he2022stgc}. ST-CGI jointly learns a predictive graph and autoregressive processes \cite{zhang2022stcgi}, while MSA-GCN incorporates Granger testing within a multistage graph-processing architecture \cite{feng2024msagcn}. These studies show that graph construction based on incremental predictive information is an established direction.}

\prev{Our ridge screen uses the same general predictive principle with a deliberately restricted feature set. A source is useful if adding its recent-speed summary reduces ranking-block MAE beyond a target's own history, neighbourhood speed and time of day. We do not perform a formal Granger significance test or infer physical causation. The comparison instead asks whether a resistance-based shortlist improves link selection when the downstream predictive screen is held fixed.}

\subsection{Position of this study}
\prev{Our study combines structural and predictive link-selection criteria within a fixed forecasting protocol. The initial greedy rule weights total-resistance reduction by directional lead--lag correlation. The ridge-based comparison then changes only the shortlist criterion while retaining the final ranking procedure. Matched random links and GTR test alternative placement rules, and a global hub provides a different communication structure. Input replacement examines source reliance separately from predictive accuracy. Table~\ref{tab:literature_position} summarizes the closest methodological relationships. This positions the work as an empirical study of explicit sensor-link selection, with over-squashing as a motivation whose role must be assessed through additional evidence.}

\begin{table*}[!t]
\centering
\caption{\prev{Representative approaches and their relation to this study. Literature comparisons describe published designs; they do not indicate that every method was implemented as a baseline.}}
\label{tab:literature_position}
\footnotesize
\setlength{\tabcolsep}{4pt}
\begin{tabular}{@{}p{3.1cm}p{5.0cm}p{8.0cm}@{}}
\toprule
\prev{Approach} & \prev{Communication or selection principle} & \prev{Relationship to this study} \\
\midrule
\prev{Learned graphs \cite{wu2019graphwavenet,wu2020mtgnn,shang2021gts,bai2020agcrn}} & \prev{Infer sensor relationships during forecasting-model training} & \prev{Our links are selected in advance. Graph WaveNet with adaptive adjacency disabled is the main backbone; the other architectures provide context.} \\
\prev{GTR \cite{black2023resistance}} & \prev{Greedily reduce total effective resistance} & \prev{Implemented structural control. Our initial selector also uses directional lead--lag correlation.} \\
\prev{FoSR and LASER \cite{karhadkar2023fosr,barbero2024laser}} & \prev{Improve spectral connectivity or balance connectivity, locality and sparsity} & \prev{Context for structural rewiring. Our selectors additionally use traffic observations and target low-speed regression.} \\
\prev{DRew \cite{gutteridge2023drew}} & \prev{Layer-dependent multi-hop aggregation with delays} & \prev{Context for distance-sensitive communication. Our augmentation is fixed across layers and training.} \\
\prev{STGC-GNNs and ST-CGI \cite{he2022stgc,zhang2022stcgi}} & \prev{Construct or learn predictive dependencies using Granger-style analysis} & \prev{Closest predictive-graph context. Our ridge screen ranks a fixed budget of additions and is not a formal causal test.} \\
\prev{Virtual nodes \cite{cao2025virtual}} & \prev{Learn hub connections using adaptive or semi-adaptive adjacency} & \prev{Inspiration for the implemented fixed global-hub control; its added-edge budget exceeds that of the sparse selectors.} \\
\prev{PDFormer and STAEformer \cite{jiang2023pdformer,liu2023staeformer}} & \prev{Use attention, propagation-delay modelling or adaptive embeddings} & \prev{Context for alternative traffic architectures. Our study keeps the backbone unchanged within each graph comparison.} \\
\prev{This study} & \prev{Combine resistance with lead--lag evidence, or compare resistance and predictive shortlists} & \prev{Evaluate frozen directed additions, architecture transfer, source-reliance diagnostics and validation-based tuning.} \\
\bottomrule
\end{tabular}
\end{table*}
